\documentclass[aspectratio=169,10pt]{beamer}
\usepackage[utf8]{inputenc}
\usepackage[T1]{fontenc}
\usepackage{lmodern}
\usepackage{amsmath,amssymb,mathtools}
\usepackage{graphicx}
\usepackage{booktabs}
\usepackage{array}
\usepackage{tikz}
\usepackage{pgfplots}
\pgfplotsset{compat=1.18}
\usepackage{ragged2e}
\usepackage{multicol}

\definecolor{iteraBlue}{RGB}{20,76,123}
\definecolor{iteraLav}{RGB}{226,226,247}
\definecolor{iteraGreen}{RGB}{24,122,48}
\definecolor{iteraGreenLight}{RGB}{231,245,233}
\definecolor{iteraRed}{RGB}{180,35,35}
\definecolor{iteraGold}{RGB}{197,159,70}
\definecolor{iteraGray}{RGB}{90,90,90}

\setbeamertemplate{navigation symbols}{}
\setbeamertemplate{frametitle}{%
  \nointerlineskip
  \begin{beamercolorbox}[wd=\paperwidth,ht=0.72cm,dp=0.18cm,leftskip=0.45cm]{frametitle}
    \usebeamerfont{frametitle}\insertframetitle
  \end{beamercolorbox}}
\setbeamercolor{frametitle}{bg=iteraLav,fg=iteraBlue}
\setbeamerfont{frametitle}{series=\bfseries,size=\large}
\setbeamercolor{structure}{fg=iteraBlue}
\setbeamercolor{block title}{bg=iteraBlue,fg=white}
\setbeamercolor{block body}{bg=iteraLav!35,fg=black}
\setbeamercolor{example title}{bg=iteraGreen,fg=white}
\setbeamercolor{example body}{bg=iteraGreenLight,fg=black}
\setbeamercolor{alerted text}{fg=iteraRed}
\setbeamertemplate{blocks}[rounded][shadow=true]
\setbeamertemplate{itemize item}{\color{iteraBlue}$\bullet$}
\setbeamertemplate{itemize subitem}{\color{iteraBlue}$\circ$}
\setbeamerfont{footline}{size=\tiny}
\setbeamertemplate{footline}{%
  \leavevmode\hbox{%
  \begin{beamercolorbox}[wd=.78\paperwidth,ht=2.6ex,dp=1.1ex,leftskip=.45cm]{author in head/foot}%
  \color{iteraGray}Matematika 3 - KL25-21001 | Kelas RB | Rifky Fauzi
  \end{beamercolorbox}%
  \begin{beamercolorbox}[wd=.22\paperwidth,ht=2.6ex,dp=1.1ex,rightskip=.45cm plus1fil]{date in head/foot}%
  \hfill\color{iteraGray}\insertframenumber/\inserttotalframenumber
  \end{beamercolorbox}}}

\title{Pertemuan 6\\Transformasi Fourier II: Strategi Deret Fourier untuk Fungsi Potongan}
\subtitle{Fokus pada gelombang kotak dan gelombang gergaji}
\author{Rifky Fauzi}
\institute{Program Studi Teknik Kelautan\\Institut Teknologi Sumatera}
\date{Kelas RB | Kamis, 14.50--17.30 | E302}

\newcommand{\ds}{\displaystyle}
\newcommand{\p}{\partial}

\begin{document}

\begin{frame}[plain]
\centering
\vspace{0.4cm}
\begin{beamercolorbox}[wd=.93\textwidth,sep=10pt,rounded=true,shadow=true]{frametitle}
\centering\Large\bfseries Pertemuan 6\\[2mm]
\large Transformasi Fourier II: Strategi Deret Fourier untuk Fungsi Potongan
\end{beamercolorbox}
\vspace{0.35cm}
\large Fokus: gelombang kotak dan gelombang gergaji
\vspace{0.55cm}

\normalsize Rifky Fauzi\\[1mm]
Program Studi Teknik Kelautan\\Institut Teknologi Sumatera
\vspace{0.35cm}

\textbf{Kelas RB} \quad | \quad Kamis, 14.50--17.30 \quad | \quad E302
\end{frame}

\begin{frame}{Fokus Pertemuan Hari Ini}
\begin{block}{Target}
Setelah pertemuan ini, mahasiswa diharapkan lebih terarah ketika mengerjakan soal deret Fourier untuk fungsi potongan, khususnya:
\begin{itemize}
  \item gelombang kotak dengan periode $2\pi$,
  \item variasi gelombang kotak dengan dua level berbeda,
  \item gelombang gergaji.
\end{itemize}
\end{block}
\begin{alertblock}{Strategi utama}
Jangan langsung mengintegralkan. Mulai dari: \textbf{(1) tentukan satu periode}, \textbf{(2) tulis fungsi secara piecewise}, \textbf{(3) cek simetri}, \textbf{(4) tentukan $a_0$, $a_n$, dan $b_n$ secara bertahap}. 
\end{alertblock}
\end{frame}

\begin{frame}{Rumus Dasar Deret Fourier}
Untuk fungsi periodik berperiode $2\pi$,
\[
f(x) \sim \frac{a_0}{2}+\sum_{n=1}^{\infty}\left(a_n\cos nx+b_n\sin nx\right),
\]
dengan
\[
a_0=\frac{1}{\pi}\int_0^{2\pi}f(x)\,dx,
\quad
 a_n=\frac{1}{\pi}\int_0^{2\pi}f(x)\cos(nx)\,dx,
\quad
 b_n=\frac{1}{\pi}\int_0^{2\pi}f(x)\sin(nx)\,dx.
\]
\begin{block}{Catatan}
Hari ini interval kerja utama kita dibuat pada $0<x<2\pi$, lalu diperluas periodik.
\end{block}
\end{frame}

\begin{frame}{Langkah Kerja yang Konsisten}
\begin{enumerate}
  \item Gambar satu periode dengan jelas.
  \item Tulis fungsi piecewise pada interval kerja.
  \item Putuskan apakah ada simetri yang bisa membantu.
  \item Hitung $a_0$ dulu untuk menangkap nilai rata-rata.
  \item Hitung $a_n$ dan $b_n$ dengan membagi integral sesuai potongan fungsi.
  \item Rapikan hasil, lalu tafsirkan pola harmoniknya.
\end{enumerate}
\begin{exampleblock}{Poin penting}
Pada fungsi potongan, kesalahan paling sering terjadi bukan pada rumus, tetapi pada \textbf{batas integral} dan \textbf{penulisan fungsi pada tiap subinterval}.
\end{exampleblock}
\end{frame}

\begin{frame}{Contoh 1: Gelombang Kotak Sederhana}
Tinjau fungsi periodik berperiode $2\pi$:
\[
f(x)=
\begin{cases}
0, & 0<x<\pi,\\
H, & \pi<x<2\pi,
\end{cases}
\qquad f(x+2\pi)=f(x).
\]
\begin{center}
\begin{tikzpicture}[scale=0.82]
\draw[->] (-0.5,0) -- (7.1,0) node[right] {$x$};
\draw[->] (0,-0.4) -- (0,2.0) node[above] {$f(x)$};
\draw[thick] (0,0) -- (3.1416,0);
\draw[thick] (3.1416,1.2) -- (6.2832,1.2);
\draw[dashed] (0,0)--(0,1.2);
\draw[dashed] (3.1416,0)--(3.1416,1.2);
\draw[dashed] (6.2832,0)--(6.2832,1.2);
\node at (4.7,1.45) {$H$};
\node[below] at (0,0) {$0$};
\node[below] at (3.1416,0) {$\pi$};
\node[below] at (6.2832,0) {$2\pi$};
\end{tikzpicture}
\end{center}
\end{frame}

\begin{frame}{Strategi untuk Contoh 1}
Untuk fungsi
\[
f(x)=
\begin{cases}
0, & 0<x<\pi,\\
H, & \pi<x<2\pi,
\end{cases}
\]
langkah kerja yang paling aman adalah:
\begin{enumerate}
  \item hitung \(a_0\) dari luas total dalam satu periode,
  \item bagi integral \(a_n\) menjadi dua bagian: \((0,\pi)\) dan \((\pi,2\pi)\),
  \item lakukan hal yang sama untuk \(b_n\),
  \item amati apakah suku tertentu hilang karena faktor \(\sin n\pi\) atau \(\cos n\pi\).
\end{enumerate}
\begin{block}{Latihan papan}
Sebelum melihat hasil, cobalah tentukan sendiri:
\[
a_0=\ ? \qquad a_n=\ ? \qquad b_n=\ ?
\]
\end{block}
\end{frame}


\begin{frame}{Perhitungan Lengkap Contoh 1}
\scriptsize
Untuk
\[
f(x)=\begin{cases}0,&0<x<\pi,\\ H,&\pi<x<2\pi,
\end{cases}
\qquad f(x+2\pi)=f(x).
\]
\textbf{Koefisien konstan}
\[
a_0=\frac{1}{\pi}\int_0^{2\pi}f(x)\,dx
=\frac{1}{\pi}\left(\int_0^\pi 0\,dx+\int_\pi^{2\pi}H\,dx\right)=H.
\]
\textbf{Koefisien kosinus}
\[
a_n=\frac{1}{\pi}\int_0^{2\pi}f(x)\cos nx\,dx
=\frac{H}{\pi}\int_\pi^{2\pi}\cos nx\,dx
=\frac{H}{n\pi}\bigl[\sin nx\bigr]_\pi^{2\pi}=0.
\]
\textbf{Koefisien sinus}
\[
b_n=\frac{H}{\pi}\int_\pi^{2\pi}\sin nx\,dx
=\frac{H}{\pi}\left[-\frac{\cos nx}{n}\right]_\pi^{2\pi}
=\frac{H}{n\pi}\bigl(\cos n\pi-\cos 2n\pi\bigr).
\]
Karena \(\cos n\pi=(-1)^n\) dan \(\cos 2n\pi=1\), maka
\[
b_n=\frac{H}{n\pi}\bigl((-1)^n-1\bigr).
\]
\end{frame}

\begin{frame}{Hasil Contoh 1}
Diperoleh
\[
a_0=H,
\qquad
 a_n=0,
\qquad
 b_n=\frac{H}{n\pi}\bigl((-1)^n-1\bigr).
\]
Jadi
\[
f(x)\sim \frac{H}{2}+\sum_{n=1}^{\infty}\frac{H}{n\pi}\bigl((-1)^n-1\bigr)\sin(nx).
\]
Karena
\[
(-1)^n-1=
\begin{cases}
0, & n \text{ genap},\\
-2, & n \text{ ganjil},
\end{cases}
\]
seri dapat ditulis sebagai
\[
f(x)\sim \frac{H}{2}-\frac{2H}{\pi}\left(\sin x+\frac{1}{3}\sin 3x+\frac{1}{5}\sin 5x+\cdots\right).
\]
\end{frame}

\begin{frame}{Contoh 2: Dua Level Berbeda}
Sekarang pertimbangkan fungsi periodik berperiode $2\pi$:
\[
f(x)=
\begin{cases}
2, & 0<x<\pi,\\
-1, & \pi<x<2\pi,
\end{cases}
\qquad f(x+2\pi)=f(x).
\]
\begin{center}
\begin{tikzpicture}[scale=0.82]
\draw[->] (-0.5,0) -- (7.1,0) node[right] {$x$};
\draw[->] (0,-1.6) -- (0,2.2) node[above] {$f(x)$};
\draw[thick] (0,1.3) -- (3.1416,1.3);
\draw[thick] (3.1416,-0.8) -- (6.2832,-0.8);
\draw[dashed] (0,0)--(0,1.3);
\draw[dashed] (3.1416,-0.8)--(3.1416,1.3);
\draw[dashed] (6.2832,0)--(6.2832,-0.8);
\node at (1.57,1.58) {$2$};
\node at (4.7,-1.08) {$-1$};
\node[below] at (0,0) {$0$};
\node[below] at (3.1416,0) {$\pi$};
\node[below] at (6.2832,0) {$2\pi$};
\end{tikzpicture}
\end{center}
\end{frame}

\begin{frame}{Strategi untuk Contoh 2}
\begin{block}{Apa yang berubah?}
Karena level bawah sekarang bernilai negatif, rata-rata fungsi ikut berubah. Artinya, suku konstan \(a_0/2\) harus benar-benar diperhatikan.
\end{block}
\begin{enumerate}
  \item Hitung \(a_0\) dari jumlah luas bertanda.
  \item Hitung \(a_n\): biasanya masih sederhana karena muncul integral kosinus pada dua potongan.
  \item Hitung \(b_n\): di sini biasanya pola harmonik lebih terlihat.
\end{enumerate}
\begin{alertblock}{Arah berpikir}
Kalau level bawah tidak nol, jangan buru-buru berharap hasilnya sama persis dengan contoh sebelumnya. \(a_0\) hampir pasti berubah.
\end{alertblock}
\end{frame}


\begin{frame}{Perhitungan Lengkap Contoh 2}
\scriptsize
Untuk
\[
f(x)=\begin{cases}2,&0<x<\pi,\\ -1,&\pi<x<2\pi,
\end{cases}
\qquad f(x+2\pi)=f(x).
\]
\textbf{Koefisien konstan}
\[
a_0=\frac{1}{\pi}\left(\int_0^\pi 2\,dx+\int_\pi^{2\pi}(-1)\,dx\right)=\frac{1}{\pi}(2\pi-\pi)=1.
\]
\textbf{Koefisien kosinus}
\[
a_n=\frac{1}{\pi}\left(2\int_0^\pi\cos nx\,dx-\int_\pi^{2\pi}\cos nx\,dx\right)
\]
\[
=\frac{1}{n\pi}\left(2[\sin nx]_0^\pi-[\sin nx]_\pi^{2\pi}\right)=0.
\]
\textbf{Koefisien sinus}
\[
b_n=\frac{1}{\pi}\left(2\int_0^\pi\sin nx\,dx-\int_\pi^{2\pi}\sin nx\,dx\right)
\]
\[
=\frac{1}{\pi}\left(\frac{2(1-(-1)^n)}{n}-\frac{(-1)^n-1}{n}\right)
=\frac{3}{n\pi}\bigl(1-(-1)^n\bigr).
\]
\end{frame}

\begin{frame}{Hasil Contoh 2}
Untuk
\[
f(x)=
\begin{cases}
2, & 0<x<\pi,\\
-1, & \pi<x<2\pi,
\end{cases}
\]
diperoleh
\[
a_0=1,
\qquad a_n=0,
\qquad b_n=\frac{3}{n\pi}\bigl(1-(-1)^n\bigr).
\]
Maka
\[
f(x)\sim \frac{1}{2}+\sum_{n=1}^{\infty}\frac{3}{n\pi}\bigl(1-(-1)^n\bigr)\sin(nx)
\]
atau
\[
f(x)\sim \frac{1}{2}+\frac{6}{\pi}\left(\sin x+\frac{1}{3}\sin 3x+\frac{1}{5}\sin 5x+\cdots\right).
\]
\end{frame}

\begin{frame}{Dari Grafik ke Fungsi Piecewise}
Kalau yang diberikan adalah grafik seperti pada soal, biasakan urutan berikut:
\begin{enumerate}
  \item Tentukan satu periode kerja, misalnya \((0,2\pi)\).
  \item Baca tinggi tiap ruas secara hati-hati.
  \item Tentukan titik loncatan: di mana fungsi berubah nilai?
  \item Tulis fungsi piecewise dengan lengkap, lalu baru masuk ke rumus Fourier.
\end{enumerate}
\begin{exampleblock}{Kebiasaan yang bagus}
Sebelum menghitung koefisien, tuliskan dulu: \textit{``Pada interval mana fungsi bernilai apa?''}
\end{exampleblock}
\end{frame}

\begin{frame}{Contoh 3: Gelombang Gergaji}
Sekarang lihat fungsi periodik berperiode $2\pi$:
\[
f(x)=x,
\qquad 0<x<2\pi,
\qquad f(x+2\pi)=f(x).
\]
\begin{center}
\begin{tikzpicture}[scale=0.82]
\draw[->] (-0.5,0) -- (7.1,0) node[right] {$x$};
\draw[->] (0,-0.3) -- (0,2.6) node[above] {$f(x)$};
\draw[thick] (0,0) -- (6.2832,2.1);
\draw[dashed] (6.2832,0)--(6.2832,2.1);
\node[below] at (0,0) {$0$};
\node[below] at (3.1416,0) {$\pi$};
\node[below] at (6.2832,0) {$2\pi$};
\node[left] at (0,2.1) {$2\pi$};
\end{tikzpicture}
\end{center}
\end{frame}

\begin{frame}{Strategi untuk Gelombang Gergaji}
Pada fungsi ini, kita tidak lagi berhadapan dengan fungsi bertingkat, tetapi prinsipnya sama:
\begin{enumerate}
  \item tentukan satu periode,
  \item hitung \(a_0\), \(a_n\), dan \(b_n\) dari integral pada interval \((0,2\pi)\),
  \item gunakan integrasi parsial saat diperlukan.
\end{enumerate}
\begin{block}{Latihan papan}
Coba kerjakan dulu:
\[
a_0=\ ? \qquad a_n=\ ? \qquad b_n=\ ?
\]
setidaknya sampai tahap integral.
\end{block}
\end{frame}


\begin{frame}{Perhitungan Lengkap Gelombang Gergaji}
\scriptsize
Untuk
\[
f(x)=x,
\qquad 0<x<2\pi,
\qquad f(x+2\pi)=f(x).
\]
\textbf{Koefisien konstan}
\[
a_0=\frac{1}{\pi}\int_0^{2\pi}x\,dx
=\frac{1}{\pi}\left[\frac{x^2}{2}\right]_0^{2\pi}=2\pi.
\]
\textbf{Koefisien kosinus} dengan integrasi parsial:
\[
a_n=\frac{1}{\pi}\int_0^{2\pi}x\cos nx\,dx.
\]
Ambil \(u=x\), \(dv=\cos nx\,dx\), maka \(du=dx\), \(v=\sin nx/n\).
\[
a_n=\frac{1}{\pi}\left(\left[\frac{x\sin nx}{n}\right]_0^{2\pi}-\frac{1}{n}\int_0^{2\pi}\sin nx\,dx\right)=0.
\]
\textbf{Koefisien sinus}:
\[
b_n=\frac{1}{\pi}\int_0^{2\pi}x\sin nx\,dx.
\]
Ambil \(u=x\), \(dv=\sin nx\,dx\), maka \(du=dx\), \(v=-\cos nx/n\).
\[
b_n=\frac{1}{\pi}\left(\left[-\frac{x\cos nx}{n}\right]_0^{2\pi}+\frac{1}{n}\int_0^{2\pi}\cos nx\,dx\right)=-\frac{2}{n}.
\]
\end{frame}

\begin{frame}{Hasil Gelombang Gergaji}
Untuk
\[
f(x)=x,
\qquad 0<x<2\pi,
\qquad f(x+2\pi)=f(x),
\]
diperoleh
\[
a_0=2\pi,
\qquad a_n=0,
\qquad b_n=-\frac{2}{n}.
\]
Jadi deret Fouriernnya adalah
\[
f(x)\sim \pi-2\sum_{n=1}^{\infty}\frac{1}{n}\sin(nx).
\]
\begin{exampleblock}{Makna}
Gelombang gergaji bisa dipandang sebagai penjumlahan tak hingga dari gelombang sinus dengan amplitudo yang menurun seperti \(1/n\).
\end{exampleblock}
\end{frame}

\begin{frame}{Bagaimana Menjawab Soal Aplikasi?}
Untuk soal seperti \textit{``Jelaskan kegunaan deret Fourier dalam teknik kelautan''}, pola jawaban yang rapi bisa seperti ini:
\begin{enumerate}
  \item Deret Fourier memecah fenomena periodik menjadi komponen harmonik.
  \item Dalam teknik kelautan, ini berguna untuk menganalisis gelombang, pasang surut, dan sinyal getaran struktur.
  \item Koefisien Fourier membantu mengenali frekuensi dominan dan besar kontribusinya.
  \item Hasil analisis dapat dipakai untuk identifikasi resonansi, filtrasi sinyal, dan interpretasi data pengukuran laut.
\end{enumerate}
\end{frame}

\begin{frame}{Ringkasan Strategi}
\begin{enumerate}
  \item Kalau grafik berbentuk kotak: fokus pada penulisan fungsi piecewise dan batas integral.
  \item Kalau level bawah tidak nol: perhatikan perubahan \(a_0\).
  \item Kalau grafik berbentuk gergaji: siapkan integrasi parsial untuk mencari \(b_n\).
  \item Selalu mulai dari satu periode yang jelas, misalnya \((0,2\pi)\).
  \item Setelah koefisien diperoleh, baru tulis deret Fourier lengkapnya.
\end{enumerate}
\begin{alertblock}{Untuk latihan berikutnya}
Setelah strategi ini dikuasai, pertemuan berikutnya bisa benar-benar difokuskan pada latihan dan variasi soal.
\end{alertblock}
\end{frame}

\end{document}
